\documentclass[12pt]{article}
\usepackage[T1]{fontenc}
\usepackage[utf8]{inputenc}
\usepackage{lmodern}
\usepackage{textcomp}
\usepackage{enumitem}
\usepackage{geometry}
\geometry{margin=1in}
\begin{document}
\begin{center}
Answer Key - Constitutional Law - Undergraduate Level Assessment 3 \
\small (Answers are ordered exactly as in the question paper.)
\end{center}

\section*{Multiple Choice}
\begin{enumerate}[label=\arabic*.]
\item \textbf{Answer:} B
\\ \textit{Options:} A) Historical; B) Analytical; C) Command; D) Sociological
\\ \textit{Note:} The Analytical School of jurisprudence emphasizes that legal principles and the law itself are derived from logical reasoning and systematic analysis, distinguishing it from other schools that may focus on moral or social considerations.
\item \textbf{Answer:} B
\\ \textit{Options:} A) To advance the interests of politicians.; B) To produce virtuous citizens and encourage righteousness in individuals.; C) To obtain power.; D) To secure justice by abolishing slavery.
\\ \textit{Note:} Aristotle believed that the ultimate aim of politics is to cultivate virtue among citizens, fostering a moral community where individuals pursue righteousness and the common good.
\item \textbf{Answer:} C
\\ \textit{Options:} A) Because rehabilitation of offenders presumes a recognition of the role of vengeance.; B) Because the lex talionis is misguided.; C) Because other justifications appear to have failed.; D) Because revenge is anachronistic.
\\ \textit{Note:} The increasing recognition of retribution as a fundamental element of punishment stems from the perception that alternative justifications, such as deterrence or rehabilitation, have not effectively addressed crime or reduced recidivism rates.
\item \textbf{Answer:} C
\\ \textit{Options:} A) Cause and effect; B) Theory and fact; C) Being obliged and having an obligation; D) Corporeal and incorporeal rights
\\ \textit{Note:} Hart's analysis of law distinguishes between being obliged, which refers to the actual behavior of individuals under coercive force or social pressure, and having an obligation, which signifies a moral or legal duty that exists independently of external enforcement, highlighting the difference between legal compliance and the normative reasons for following the law.
\item \textbf{Answer:} C
\\ \textit{Options:} A) The prohibition of the use of force; B) The prohibition of torture; C) The prohibition of genocide; D) The principle of self-determination
\\ \textit{Note:} The prohibition of genocide was explicitly recognized as a rule of jus cogens by the International Court of Justice in the 2007 case concerning Bosnia and Herzegovina v. Serbia and Montenegro, establishing it as a foundational principle of international law that cannot be violated.
\end{enumerate}

\section*{True / False}
\begin{enumerate}[label=\arabic*.]
\setcounter{enumi}{5}
\item \textbf{Answer:} False
\\ \textit{Options:} A) True; B) False
\\ \textit{Note:} John Locke argued that natural rights, such as life, liberty, and property, are fundamental to governance and must be protected by the state, thereby making the statement false.
\item \textbf{Answer:} True
\\ \textit{Options:} A) True; B) False
\\ \textit{Note:} Recognition is merely an acknowledgment by other states of an entity's status as a state, and it does not create or alter the legal existence of that state, which is determined by criteria such as a defined territory, permanent population, and effective government.
\item \textbf{Answer:} True
\\ \textit{Options:} A) True; B) False
\\ \textit{Note:} In Dworkin's framework, principles are understood as fundamental rights that inform legal decisions, while policies represent broader societal goals that guide legislative action, thus distinguishing the two concepts.
\item \textbf{Answer:} False
\\ \textit{Options:} A) True; B) False
\\ \textit{Note:} Leopold Pospisil identifies 'dispute' as a key element of legal systems, but it is not classified as one of the four essential elements that constitute the law itself.
\item \textbf{Answer:} False
\\ \textit{Options:} A) True; B) False
\\ \textit{Note:} Olivecrona's argument is considered false because it lacks the substantial empirical evidence necessary to support the psychological basis of law, relying more on philosophical reasoning than on concrete psychological findings.
\end{enumerate}

\section*{Long Form Response}
\begin{enumerate}[label=\arabic*.]
\setcounter{enumi}{10}
\item \textbf{Answer:} Rousseau's assertion that the "Social Contract is not a historical fact but a hypothetical construction of reason" suggests that the legitimacy of political authority is derived not from historical precedents but from rational agreement among individuals to form a collective society. This idea implies that constitutional law should be grounded in principles that reflect the general will, rather than merely historical traditions or precedents. In modern governance, this perspective encourages a focus on democratic participation and the consent of the governed, as it positions citizens as active participants in shaping the laws that govern them, rather than passive subjects of historical authority. Thus, Rousseau's framework invites ongoing dialogue about justice, equality, and the moral foundations of political institutions, reinforcing the need for laws and constitutions to evolve in response to the collective rationality of the populace.
\\ \textit{Note:} correct because it emphasizes that Rousseau's concept of the social contract prioritizes rational consensus over historical legitimacy, thereby advocating for constitutional law that embodies the general will and influences modern governance by challenging reliance on tradition.
\item \textbf{Answer:} Allowing a reservation to the definition of torture in the International Covenant on Civil and Political Rights (ICCPR) fundamentally undermines the treaty's object and purpose, which is to safeguard human rights and dignity. Such a reservation could lead to the normalization of practices that inflict severe pain or suffering, thus permitting states to circumvent their obligations under international law. This not only jeopardizes the protection of individuals from inhumane treatment but also erodes the credibility and effectiveness of the ICCPR as a framework for promoting human rights globally. In contemporary legal practice, maintaining a strict definition of torture is essential for upholding the principles of justice and accountability, making any reservation unacceptable.
\\ \textit{Note:} correct because allowing a reservation to the definition of torture in the ICCPR contradicts the treaty's fundamental goal of protecting human rights, potentially enabling states to justify inhumane treatment and evade their international legal responsibilities.
\end{enumerate}

\vfill
\noindent\textit{End of Answer Key}
\end{document}